\documentclass[text06.tex]{subfiles}
\begin{document}
\subsection{Juliusが起動しなくなったら}
もしもJuliusがうまく起動しなくなったときは、インストールしてあるコマンド\\
\code{stopjulius.sh}\\
を一度実行してからもう一度起動を試してみましょう。このコマンドは、上手く終了できなかった前のJuliusを終了させるコマンドです。

\subsection{HSPでのファイルの指定}
このプログラムでは、\\
\code{init\_julius "kudamono.dic"}\quad\\
として、プログラムを実行しているディレクトリにある\code{kudamono.dic}ファイルを辞書として指定しています。dicファイルが格納されているディレクトリが、プログラムを実行しているディレクトリと異なる場合は、相対パスか絶対パスでディレクトリを指定します。HSPプログラムからファイルを指定する際に\textasciitilde は使用できないので注意しましょう。このプログラムで読み込んでいるkudamono.dicを絶対パスで指定すると、\\
\code{init\_julius "/home/<ユーザ名>/07/kudamono.dic"}\\
となります。


\end{document}
